\addtocontents{toc}{\protect\newpage}

\section{Presupuesto} \label{sec:5_presupuesto}
Se detalla en este apartado la estimación del presupuesto derivado de la realización de este Trabajo Fin de Máster.
Las partidas presupuestarias se han agrupado en tres capítulos: hardware, software y mano de obra. Los precios reflejados 
no incluyen IVA.

\subsection{Capítulo 1: Hardware} \label{sec:5_presupuesto_hardware}
En este capítulo se incluyen los costes derivados de la adquisición de los equipos informáticos necesarios 
para el desarrollo del TFM, así como los periféricos asociados. Se muestran en la \autoref{tab:hardware_budget}.

\begin{table}[h!]
	\centering
	\renewcommand{\arraystretch}{1.35}
	\renewcommand{\tabularxcolumn}[1]{m{#1}}
	\setlength{\tabcolsep}{8pt}
	\begin{tabularx}{\linewidth}{|c|>{\hsize=1.785\hsize\centering\arraybackslash}X|c|c|>{\hsize=0.6\hsize\centering\arraybackslash}X|>{\hsize=0.615\hsize\centering\arraybackslash}X|}
		\hline
		\rowcolor[gray]{0.90}
		\textbf{Id}                          & \textbf{Descripci\'on}                           & \textbf{Unidad} & \textbf{Cantidad} & \textbf{Precio unitario} & \textbf{Total}       \\
		\hline
		\rowcolor{white}
		1                                    & ASUS Vivobook M6500QC                            & u               & 1                 & 1.249,99 \texteuro{}     & 1.249,99 \texteuro{} \\
		\hline
		\rowcolor{white}
		1.1                                  & Procesador AMD Ryzen 7 5800 H                    & u               & 1                 & -                        & -                    \\
		\rowcolor{white}
		1.2                                  & Tarjeta gr\'afica NVIDIA GeForce RTX 3050 Mobile & u               & 1                 & -                        & -                    \\
		\rowcolor{white}
		1.3                                  & Pantalla 15,6'' FHD (1920 x 1080) OLED           & u               & 1                 & -                        & -                    \\
		\rowcolor{white}
		1.4                                  & 16 GB RAM DDR4                                   & u               & 1                 & -                        & -                    \\
		\rowcolor{white}
		1.5                                  & SSD 512 GB                                       & u               & 1                 & -                        & -                    \\
		\rowcolor{white}
		1.6                                  & Teclado integrado                                & u               & 1                 & -                        & -                    \\
		\hline
		\rowcolor{white}
		2                                    & Monitor MSI MAG 275CQF WQHD 27'' 180Hz           & u               & 1                 & 229,00 \texteuro{}       & 229,00 \texteuro{}   \\
		\hline
		\rowcolor{white}
		3                                    & Teclado mec\'anico Keychron K2                   & u               & 1                 & 106,99 \texteuro{}       & 106,99 \texteuro{}   \\
		\hline
		\rowcolor{white}
		4                                    & Rat\'on inal\'ambrico Logitech MX Master 3S      & u               & 1                 & 99,99 \texteuro{}        & 99,99 \texteuro{}    \\
		\hline
		\rowcolor[gray]{0.96}
		\multicolumn{5}{|c|}{\textbf{TOTAL}} & \textbf{1.685,97 \texteuro{}}                                                                                                            \\
		\hline
	\end{tabularx}
	\caption{Capítulo presupuestario 1: Hardware}
	\label{tab:hardware_budget}
\end{table}


El coste total de este capítulo equivale a MIL SEISCIENTOS OCHENTA Y CINCO EUROS CON NOVENTA Y SIETE CÉNTIMOS.

\newpage

\subsection{Capítulo 2: Software} \label{sec:5_presupuesto_software}
En este capítulo se incluyen los costes derivados de la adquisición de las licencias de software necesarias
para el desarrollo del TFM, así como cualquier otro software requerido en la implementación y desarrollo. No 
se incluyen licencias de software libre o de código abierto. Se muestran en la \autoref{tab:software_budget}.

\begin{table}[h!]
	\centering
	\renewcommand{\arraystretch}{1.35}
	\renewcommand{\tabularxcolumn}[1]{m{#1}}
	\setlength{\tabcolsep}{8pt}
	\begin{tabularx}{\linewidth}{|c|>{\hsize=1.785\hsize\centering\arraybackslash}X|c|c|>{\hsize=0.6\hsize\centering\arraybackslash}X|>{\hsize=0.615\hsize\centering\arraybackslash}X|}
		\hline
		\rowcolor[gray]{0.90}
		\textbf{Id}                          & \textbf{Descripci\'on}                       & \textbf{Unidad} & \textbf{Cantidad} & \textbf{Precio unitario} & \textbf{Total}     \\
		\hline
		\rowcolor{white}
		1                                    & Licencia Microsoft Office 365 Personal       & u / a\~no       & 1                 & 99,00 \texteuro{}        & 99,00 \texteuro{}  \\
		\hline
		\rowcolor{white}
		2                                    & GitHub Student Developer Pack                & u               & 1                 & 0,00 \texteuro{}         & 0,00\texteuro{}    \\
		\hline
		\rowcolor[gray]{0.96}
		\multicolumn{5}{|c|}{\textbf{TOTAL}} & \textbf{99,00 \texteuro{}}                                                                                                        \\
		\hline
	\end{tabularx}
	\caption{Capítulo presupuestario 2: Software}
	\label{tab:software_budget}
\end{table}


El coste total de este capítulo equivale a NOVENTA Y NUEVE EUROS.

\subsection{Amortización del inmovilizado material} \label{sec:5_presupuesto_amortizacion}
En los capítulos anteriores se han detallado los precios asociados a la adquisición de los recursos de hardware y 
software necesarios para el desarrollo del TFM. Sin embargo, en lugar de considerar el coste total de estos recursos 
como gasto único, se ha optado por calcular su coste de amortización.

Se ha determinado un periodo de amortización de seis años. De acuerdo con la planificación 
temporal del proyecto (\autoref{sec:4_planificacion}), el tiempo de utilización efectiva de estos recursos 
se estima en cuatro meses. Esto coincide con la duración total del proyecto, permitiendo así una distribución más
precisa de los costes y reflejando el uso real de los recursos a lo largo del mismo. En la \autoref{tab:depreciation_budget} 
se muestra el cálculo detallado de la amortización del inmovilizado material.

\begin{table}[h!]
	\centering
	\renewcommand{\arraystretch}{1.35}
	\renewcommand{\tabularxcolumn}[1]{m{#1}}
	\setlength{\tabcolsep}{8pt}
	\begin{tabularx}{\linewidth}{|>{\hsize=0.9\hsize\centering\arraybackslash}X|>{\hsize=1\hsize\centering\arraybackslash}X|>{\hsize=1.1\hsize\centering\arraybackslash}X|>{\hsize=1\hsize\centering\arraybackslash}X|>{\hsize=1\hsize\centering\arraybackslash}X|}
		\hline
		\rowcolor[gray]{0.90}
		\textbf{Material}                                   & \textbf{Cuota de adquisici\'on} & \textbf{Tiempo de amortizaci\'on (a\~nos)} & \textbf{Tiempo de uso (meses)} & \textbf{Amortizaci\'on} \\
		\hline
		\rowcolor{white}
		Hardware                                            & 1.685,97 \texteuro{}            & 6                                          & 4                              & 93,67  \texteuro{}      \\
		\hline
		\rowcolor{white}
		Software                                            & 99,00 \texteuro{}               & 6                                          & 4                              & 5,50  \texteuro{}       \\
		\hline
		\rowcolor[gray]{0.96}
		\multicolumn{4}{|c|}{\textbf{AMORTIZACI\'ON TOTAL}} & \textbf{99,17 \texteuro{}}                                                                                                              \\
		\hline
	\end{tabularx}
	\caption{Amortización del inmovilizado material}
	\label{tab:depreciation_budget}
\end{table}


El coste total correspondiente a la amortización del inmovilizado material equivale a NOVENTA Y NUEVE EUROS CON DIECISIETE CÉNTIMOS.

\newpage

\subsection{Capítulo 3: Mano de obra} \label{sec:5_presupuesto_mano_obra}
En este capítulo se recogen los costes asociados a la mano de obra necesaria para el desarrollo del TFM. Se han considerado 
distintos perfiles profesionales, definiendo para cada uno de ellos su correspondiente tarifa horaria y una estimación del número de horas de dedicación.
Se muestran en la \autoref{tab:workforce_budget}.

\begin{table}[h!]
	\centering
	\renewcommand{\arraystretch}{1.35}
	\renewcommand{\tabularxcolumn}[1]{m{#1}}
	\setlength{\tabcolsep}{8pt}
	\begin{tabularx}{\linewidth}{|c|>{\hsize=1.9\hsize\centering\arraybackslash}X|c|c|>{\hsize=0.425\hsize\centering\arraybackslash}X|>{\hsize=0.675\hsize\centering\arraybackslash}X|}
		\hline
		\rowcolor[gray]{0.90}
		\textbf{Id}                          & \textbf{Descripci\'on}             & \textbf{Unidad} & \textbf{Cantidad} & \textbf{Precio unitario} & \textbf{Total}       \\
		\hline
		\rowcolor{white}
		1                                    & Especialista en PyTorch/RAW Images & Horas           & 380               & 26,00 \texteuro{}        & 9.880,00 \texteuro{} \\
		\hline
		\rowcolor{white}
		2                                    & Desarrollador Python/PyQt          & Horas           & 220               & 20,00 \texteuro{}        & 4.400,00 \texteuro{} \\
		\hline
		\rowcolor{white}
		3                                    & Product Owner                      & Horas           & 40                & 30,00 \texteuro{}        & 1.200,00 \texteuro{} \\
		\hline
		\rowcolor{white}
		4                                    & Scrum master                       & Horas           & 40                & 30,00 \texteuro{}        & 1.200,00 \texteuro{} \\
		\hline
		\rowcolor[gray]{0.96}
		\multicolumn{5}{|c|}{\textbf{TOTAL}} & \textbf{16.680,00 \texteuro{}}                                                                                             \\
		\hline
	\end{tabularx}
	\caption{Capítulo presupuestario 3: Mano de obra}
	\label{tab:workforce_budget}
\end{table}


El coste total de este capítulo equivale a DIECISÉIS MIL SEISCIENTOS OCHENTA EUROS.

\subsection{Presupuesto total} \label{sec:5_presupuesto_total}
Se presenta en la \autoref{tab:total_budget} el presupuesto total del TFM. El cálculo del
presupuesto incluye los costes de amortización de hardware y software y mano de obra. Se
consideran también los gastos generales y el beneficio industrial, que se aplican como
porcentajes sobre el presupuesto de ejecución material. Los gastos generales se han estimado en
un 13\% para cubrir costes indirectos y el beneficio industrial se establece en un 6\% como margen
para el contratista por el proyecto realizado.

\begin{table}[h!]
	\centering
	\renewcommand{\arraystretch}{1.35}
	\renewcommand{\tabularxcolumn}[1]{m{#1}}
	\setlength{\tabcolsep}{8pt}
	\begin{tabularx}{\linewidth}{|>{\hsize=1.7\hsize\centering\arraybackslash}X|>{\hsize=0.65\hsize\centering\arraybackslash}X|>{\hsize=0.65\hsize\centering\arraybackslash}X|}
		\hline
		\rowcolor[gray]{0.90}
		\textbf{Cap\'itulos}                    & \textbf{Porcentaje} & \textbf{Total}                 \\
		\hline
		\rowcolor{white}
		Total inmovilizado material             & -                   & 99,17 \texteuro{}              \\
		\hline
		\rowcolor{white}
		Cap\'itulo 3: Mano de obra              & -                   & 16.680,00 \texteuro{}          \\
		\specialrule{1.1pt}{0pt}{0pt}
		\rowcolor{white}
		Presupuesto de ejecuci\'on material     & -                   & 16.779,17 \texteuro{}          \\
		\hline
		\rowcolor{white}
		Gastos generales                        & 13\%                & 2.181,29 \texteuro{}           \\
		\hline
		\rowcolor{white}
		Beneficio industrial                    & 6\%                 & 1.006,75 \texteuro{}           \\
		\specialrule{1.1pt}{0pt}{0pt}
		\rowcolor{white}
		Presupuesto de ejecuci\'on por contrata & -                   & 19.967,21 \texteuro{}          \\
		\hline
		\rowcolor{white}
		IVA                                     & 21\%                & 4.193,11 \texteuro{}           \\
		\hline
		\rowcolor[gray]{0.96}
		\textbf{PRESUPUESTO TOTAL}              &                     & \textbf{24.160,32 \texteuro{}} \\
		\hline
	\end{tabularx}
	\caption{Presupuesto total}
	\label{tab:total_budget}
\end{table}


El presupuesto total del proyecto equivale a VEINTICUATRO MIL CIENTO SESENTA EUROS CON TREINTA Y DOS CÉNTIMOS.
